\section{Visible zero coefficients and infinite oscillation}
\label{sec:oscillation}

\subsection{Coefficient visibility}

From \eqref{eq:L-product}, the \(n\)-th factor vanishes precisely at
\[
w=\pi i\ell\,2^{n+1},\qquad \ell\in\Z\setminus\{0\}.
\]
The locally uniform nonvanishing of the tail away from factor zeros proves
\eqref{eq:L-zero-set}.  Thus every off-critical-line zero is visible.  A
critical-line zero is invisible only if its ordinate is in
\[
4\pi\Z\setminus\{0\}.
\]

Conrey's Theorem~1 in \cite{Conrey1989} proves that a positive proportion
of all nontrivial zeros are simple and lie on the critical line.  Hence
there are \(\gg T\log T\) distinct simple critical-line zeros below height
\(T\).  The invisible lattice has only \(O(T)\) ordinates in that range,
and a simple zero occupies at most one such ordinate.  It follows
unconditionally that infinitely many critical-line zeros are visible for
the frozen packet.

\subsection{A Laplace-transform Landau lemma}

\begin{lemma}[Landau]
\label{lem:landau}
Let \(T_0\ge0\), and let
\(f:[T_0,\infty)\to[0,\infty)\) be locally integrable.  Suppose its Laplace
integral converges on a nonempty right half-line, and define
\[
F(s)=\int_{T_0}^{\infty}f(t)e^{-st}\,dt,
\]
\[
\sigma_c=
\inf\left\{
\sigma\in\R:
\int_{T_0}^{\infty}f(t)e^{-\sigma t}\,dt<\infty
\right\}.
\]
If \(\sigma_c\) is finite, then \(F\) is singular at the real point
\(s=\sigma_c\).
\end{lemma}

\begin{proof}
Assume that \(F\) is holomorphic near \(\sigma_c\).  Choose
\(\sigma_1>\sigma_c\) and \(r>0\) so that
\(\sigma_1-r<\sigma_c\) while the Taylor disk from \(\sigma_1\) to
\(\sigma_1-r\) remains inside the union of this neighborhood and
\(\Re s>\sigma_c\).  For a small \(\varepsilon>0\) with
\(\sigma_1-\varepsilon>\sigma_c\), the inequality
\(t^n\le C_{n,\varepsilon}e^{\varepsilon t}\) justifies differentiation:
\[
(-1)^nF^{(n)}(\sigma_1)
=\int_{T_0}^{\infty}t^nf(t)e^{-\sigma_1t}\,dt\ge0.
\]
Taylor's theorem and monotone convergence give
\begin{align*}
F(\sigma_1-r)
&=
\sum_{n=0}^{\infty}
\frac{(-r)^nF^{(n)}(\sigma_1)}{n!}\\
&=
\int_{T_0}^{\infty}
f(t)e^{-\sigma_1t}
\sum_{n=0}^{\infty}\frac{(rt)^n}{n!}\,dt\\
&=
\int_{T_0}^{\infty}
f(t)e^{-(\sigma_1-r)t}\,dt<\infty,
\end{align*}
contradicting the definition of \(\sigma_c\).
\end{proof}

\subsection{Oscillation of the arithmetic interaction}

The direct prime-power window gives the elementary bound
\[
|\Ih(t/2)|=O_h((1+t)e^{t/2}),\qquad t>1.
\]
Indeed, one may use
\(\sum_{n\le X}\Lambda(n)\le X\log X\).  Thus the tail transform has a
finite abscissa of convergence no greater than \(1/2\).

At each visible critical-line zero \(\rho=1/2+i\gamma\), the meromorphic
continuation \eqref{eq:B-transform} has a genuine pole at \(s=i\gamma\).
Suppose \(\Ih(t/2)\) were eventually nonnegative.  Discarding the finite
initial interval changes its transform by an entire function.  The visible
pole rules out an abscissa below \(0\), while absolute convergence gives
\[
0\le\sigma_c\le\frac12.
\]
But \eqref{eq:B-transform} is holomorphic at every real point of this
interval: the pole at \(1/2\) cancels, and
\(\zeta(u)\ne0\) for real \(1/2\le u<1\).  Lemma~\ref{lem:landau} gives a
contradiction.  Applying the same argument to \(-\Ih(t/2)\) excludes
eventual nonpositivity.

\subsection{Oscillation of the complete splitting}

Put \(\mathcal Q_h(t)=\qform(h_{t/2}^{\mathrm R},
h_{t/2}^{\mathrm L})\).  The zero-only formula gives
\[
\mathcal Q_h(t)=
\sum_\rho a_\rho e^{(\rho-1/2)t},
\qquad
a_\rho=m_\rho\Lh\left(\rho-\frac12\right).
\]
For every \(T_0>1\),
\[
\int_{T_0}^{\infty}\mathcal Q_h(t)e^{-st}\,dt
=
\sum_\rho
\frac{
a_\rho e^{(\rho-1/2-s)T_0}
}{
s-(\rho-1/2)
}.
\]
Rapid coefficient decay makes the right side normally convergent away
from its displayed poles.  It has genuine nonreal poles at every visible
critical zero and no real poles.  The same Landau argument, applied
separately to \(\mathcal Q_h\) and to \(-\mathcal Q_h\), excludes either
eventual sign.  Equation~\eqref{eq:polarization} transfers the conclusion to
\(\Delta_h\).

Both \(\Ih\) and \(\Delta_h\) are continuous.  Alternating positive and
negative points can be chosen beyond increasing bounds, producing
infinitely many disjoint sign-crossing intervals.  This does not assert
that the intervening zeros are simple or transverse.
